\documentclass[a4paper,12pt]{ctexart}
\usepackage{xcolor}
\usepackage[top=1.8cm, bottom=1.8cm, left=1.8cm, right=1.8cm]{geometry}
\usepackage{array}
\usepackage{graphicx}
\linespread{1.35}

\newcommand{\sidebox}[2]{%
  \noindent
  \begin{minipage}[t]{0.48\textwidth}
  \colorbox{blue!70!black}{\textcolor{white}{\normalsize\textbf{\hspace{0.3cm}原文\hspace{0.3cm}}}}
  \vspace{0.1cm}
  \par #1
  \end{minipage}
  \hfill
  \begin{minipage}[t]{0.48\textwidth}
  \colorbox{green!60!black}{\textcolor{white}{\normalsize\textbf{\hspace{0.3cm}译文\hspace{0.3cm}}}}
  \vspace{0.1cm}
  \par #2
  \end{minipage}
  \par
}

\newcommand{\thinktitle}{%
  \vspace{0.3cm}
  {\large\textcolor{orange!80!black}{\textbf{想一想，说一说：}}}
}

\newcommand{\writetitle}{%
  \vspace{0.3cm}
  {\large\textcolor{blue!70!black}{\textbf{写一写：}}}
}

\newcommand{\genrebox}[1]{%
  \vspace{0.2cm}
  \noindent\colorbox{purple!60!white}{\textcolor{white}{\small\textbf{\hspace{0.3cm}文体小知识\hspace{0.3cm}}}}
  \par #1
}

\newcommand{\inferbox}[1]{%
  \vspace{0.2cm}
  \noindent\colorbox{orange!40!white}{\textcolor{black!80!white}{\small\textbf{\hspace{0.3cm}细节推理\hspace{0.3cm}}}}
  \par #1
}

\title{\textcolor{blue!70!black}{\Huge\textbf{古代神话与传说}}}
\author{}
\date{}

\begin{document}
\maketitle
\thispagestyle{empty}

\section*{\textcolor{blue!70!black}{前言：古人如何解释世界？}}

同学们，在没有科学知识的古代，人们看到打雷闪电、日出日落、四季更替，心里一定充满了好奇和疑问：

\vspace{0.2cm}
\begin{quote}
\large
\textit{天为什么不会掉下来？} \\
\textit{太阳每天从哪里来，又到哪里去？} \\
\textit{人是怎么出现的？} \\
\textit{为什么会有洪水？}
\end{quote}
\vspace{0.2cm}

他们没法查百度，也没法看科普书。于是，他们用自己的想象力创造了故事来解释这一切。这些故事，就是\textbf{神话}。

神话是全人类共有的财富。每个民族都有自己的神话——中国的、希腊的、埃及的、北欧的……不同的神话里，有相似的主题，也有截然不同的想象。在这个单元里，我们将一起去认识这些古老的故事。

你可能会发现，神话虽然"假"（现实中不可能发生），但它们表达的情感却是真的——勇气、恐惧、希望、爱。这就是神话的魅力所在。

\newpage

% ============================================================
\section*{\textcolor{orange!80!black}{第一周：山海经——中国最古老的神话宝藏}}
% ============================================================

你有没有想过，中国最早的神话都藏在哪里？答案就在一部叫《山海经》的奇书里。

《山海经》可不是普通的书。它里面记载了各种各样的神兽、异人、奇山怪水——有长着九条尾巴的狐狸，有会唱歌的鱼，有只有一只脚的牛……古人相信，在遥远的地方，真的有这些神奇的存在。

今天，我们来看《山海经》里两则最有名的神话故事。

\vspace{0.5cm}

\subsection*{\textcolor{teal!80!black}{故事一：精卫填海}}

\begin{quote}
\textit{传说在很久以前，炎帝有一个小女儿，名叫女娃。有一天，女娃独自驾船去东海玩。突然，海上刮起了狂风，掀起了巨浪，小船被打翻了，女娃淹死在了大海里。}
\end{quote}

\begin{quote}
\textit{女娃死后，她的灵魂化成了一只小鸟。这只鸟长着花脑袋、白嘴巴、红脚丫，名叫"精卫"。精卫痛恨无情的大海夺去了她的生命，于是每天从西山上衔来小石子和小树枝，飞到东海上空投下去。她下定决心：总有一天，要把大海填平！}
\end{quote}

\vspace{0.2cm}
\begin{center}
\textcolor{gray}{——改编自《山海经·北山经》}
\end{center}

\vspace{0.3cm}

\sidebox{
北山经又北二百里，曰发鸠之山，其上多柘木。有鸟焉，其状如乌，文首、白喙、赤足，名曰精卫，其鸣自詨。是炎帝之少女，名曰女娃。女娃游于东海，溺而不返，故为精卫。常衔西山之木石，以堙于东海。
}{ 
再往北二百里，有一座发鸠山，山上长着很多柘树。有一种鸟生活在这里，形状像乌鸦，花脑袋、白嘴巴、红脚丫，名字叫"精卫"，它的叫声就是自己名字的读音。它本是炎帝的小女儿，名叫女娃。女娃在东海游玩时淹死了，再也没有回来，所以化成了精卫鸟。它常常衔着西山上的树枝和石子，飞去填塞东海。
}

\vspace{0.3cm}

\subsection*{\textcolor{teal!80!black}{故事二：夸父逐日}}

\begin{quote}
\textit{从前有一个巨人，名叫夸父。他有一双长腿，跑起来像风一样快。}

\textit{有一天，夸父忽然冒出一个念头："我要追上太阳，把它捉住！"于是他迈开大步，朝着太阳的方向追去。他追啊追啊，穿过高山和河流，离太阳越来越近。可是太阳像一个巨大的火球，烤得他口干舌燥。}

\textit{夸父渴极了，他跑到黄河边，一口气喝干了黄河的水，又去喝干了渭河的水。可是还不够！他继续往北方跑，想去喝大泽里的水。然而，还没有走到大泽，他就渴死了。}

\textit{夸父倒下的时候，他手里的木杖掉在了地上。说来也神奇，那根木杖落下的地方，竟然长出了一片茂密的桃林——鲜美的桃子挂在枝头，可以为后来的人解渴。}
\end{quote}

\vspace{0.2cm}
\begin{center}
\textcolor{gray}{——改编自《山海经·海外北经》}
\end{center}

\vspace{0.3cm}

\sidebox{
夸父与日逐走，入日。渴，欲得饮，饮于河、渭，河、渭不足，北饮大泽。未至，道渴而死。弃其杖，化为邓林。
}{
夸父和太阳赛跑，追到了太阳边。他口渴极了，想要喝水，于是去喝黄河和渭河的水，黄河渭河的水不够喝，他又往北想去喝大泽的水。还没有走到，就在路上渴死了。他扔下的木杖，化成了一片桃林。
}

\vspace{0.5cm}

\thinktitle
\begin{enumerate}
  \item 精卫明明是一只小鸟，大海那么大，她为什么还要去填海？
  \item 夸父为什么要追太阳？你觉得他是聪明还是傻？
  \item 这两个故事的结局有什么相似之处？
  \item 如果让你选，你更喜欢精卫还是夸父？为什么？
\end{enumerate}

\inferbox
\vspace{0.1cm}
\noindent
\textbf{#1：精卫在原文中叫"女娃"——一个普通的女孩名字。为什么要安排她是一个小女孩，而不是大人？}

\vspace{0.1cm}
\noindent
\textbf{#2：夸父的杖变成了桃林——这个结尾有什么特别的意义？试想如果夸父什么都没有留下，故事的感觉会有什么不同？}

\genrebox
\vspace{0.1cm}
\noindent
\textbf{神话的特征：}
\begin{itemize}
  \item 有超自然的力量（小鸟填海、巨人追日）
  \item 解释某种自然现象或事物的起源（为什么有桃林？）
  \item 没有作者，口口相传，最初是信以为真的
\end{itemize}

\writetitle
\vspace{0.1cm}
请用3-5句话回答：精卫在填海的时候，心里在想什么？试着用精卫的口气来写。

\newpage

% ============================================================
\section*{\textcolor{orange!80!black}{第二周：希腊神话——另一个世界的神话}}
% ============================================================

上个月我们读了中国神话，这个月我们来认识一下地球另一边的希腊人创造的神话。很有意思的是，在几千年前，相隔万里的中国人和希腊人，各自编出了自己的神话。故事不一样，但有些感觉却惊人的相似。

\vspace{0.5cm}

\subsection*{\textcolor{teal!80!black}{故事一：普罗米修斯盗火}}

\begin{quote}
\textit{在希腊神话中，人类刚被创造出来的时候，生活在一片黑暗和寒冷中。他们没有火——没有火来取暖、没有火来煮熟食物、没有火来照亮黑夜。}

\textit{有一位名叫普罗米修斯的大力神，他同情人类的处境。他知道天上的神拥有火种，但众神之王宙斯禁止把火送给人类。普罗米修斯不顾禁令，偷偷折了一根茴香枝，在太阳车经过的时候点燃了它，然后带到人间。}

\textit{人类有了火，从此过上了温暖的生活。但宙斯发现后勃然大怒。他把普罗米修斯锁在高加索山上的一根岩石柱子上，每天派一只巨鹰来啄食他的肝脏。最残忍的是，普罗米修斯的肝脏被吃掉后，第二天又会重新长出来——这样，痛苦就永远没有尽头。}

\textit{但普罗米修斯没有屈服。他在痛苦中坚持了很多很多年，直到后来英雄赫拉克勒斯路过，射死了巨鹰，才把他解救下来。}
\end{quote}

\vspace{0.3cm}

\thinktitle
\begin{enumerate}
  \item 普罗米修斯明知道会受惩罚，为什么还要把火送给人类？
  \item 你觉得宙斯惩罚普罗米修斯，是因为他"偷火"这件事本身，还是因为"违抗命令"？
  \item 对比精卫和普罗米修斯，他们有什么相同之处？
  \item 如果让你给这个故事换一个结局，你会怎么改？
\end{enumerate}

\inferbox
\vspace{0.1cm}
\noindent
\textbf{#1：文中说普罗米修斯的肝脏每天被吃掉又长出来——为什么要设计这个"重复痛苦"的设定？}

\vspace{0.1cm}
\noindent
\textbf{#2：中国神话里，受到惩罚的神/人通常怎么反应？和普罗米修斯比呢？}

\genrebox
\vspace{0.1cm}
\noindent
\textbf{比较阅读：中国神话 vs 希腊神话}
\begin{itemize}
  \item 中国神话的神大多长得像人/动物混合体（精卫是鸟），希腊神话的神更像人（有力气、有情感）
  \item 中国神话往往通过一个简短的故事讲完，希腊神话像连续剧，神与神之间有复杂的关系网
\end{itemize}

\vspace{0.5cm}

\subsection*{\textcolor{teal!80!black}{故事二：潘多拉的盒子}}

\begin{quote}
\textit{普罗米修斯盗火之后，宙斯还不解气。他决定给人类降下灾难来抵消得到火的快乐。}

\textit{他让火神造了一个美丽的女人，名叫潘多拉。在送给人类之前，每个神都送给她一样"礼物"——美丽、智慧、口才、好奇……以及一个盒子（其实是一个大罐子），并且叮嘱她：千万不能打开这个盒子！}

\textit{潘多拉来到人间后，和人类一起生活。她每天都在想：盒子里到底装了什么？为什么不能打开？她的好奇心一天比一天强烈。}

\textit{终于有一天，她忍不住了。她悄悄打开盒子——刹那间，从盒子里飞出了各种可怕的东西：疾病、痛苦、嫉妒、仇恨、谎言、衰老……所有的不幸都飞了出来，散布到人间。}

\textit{潘多拉吓坏了，赶紧关上盒子。可是已经晚了——所有的不幸都已经跑出去了。盒子里只剩下最后一样东西：希望。}

\textit{从此以后，人世间虽然有了种种苦难，但人类心中始终保留着希望。这就是宙斯无法夺走的最后一件礼物。}
\end{quote}

\vspace{0.3cm}

\thinktitle
\begin{enumerate}
  \item 潘多拉明明被告知不能打开盒子，她为什么还是打开了？
  \item 为什么盒子里最后留下的是"希望"？如果连希望也飞走了，故事会变成什么感觉？
  \item 你觉得这个故事的寓意是什么？
  \item "好奇害死猫"和潘多拉的故事是不是同一个意思？
\end{enumerate}

\inferbox
\vspace{0.1cm}
\noindent
\textbf{潘多拉的好奇心到底是好还是坏？从故事里来看，结果是坏的。但如果没有好奇心，人类会有进步吗？你怎么看？}

\writetitle
\vspace{0.1cm}
两个故事都选一个，完成以下任务（二选一）：

\noindent
\textbf{A.} 改写普罗米修斯的故事：如果你是他，你会不会去偷火？请写出你的理由（300字以内）

\noindent
\textbf{B.} 写一段"如果我是潘多拉"——当你面对那个盒子的时候，你会怎么想、怎么做？

\newpage

% ============================================================
\section*{\textcolor{orange!80!black}{第三周：诗人笔下的神话——当神话遇上诗歌}}
% ============================================================

神话原来只是口口相传的故事。后来有了文字，人们把神话记了下来。再后来，诗人们开始把这些神话写进诗里——他们不是简单地重复故事，而是借用神话来表达自己的情感。

今天我们来读两首和神话有关的诗。

\vspace{0.5cm}

\subsection*{\textcolor{teal!80!black}{嫦娥}}

\begin{center}
\large
\textbf{嫦娥} \hspace{2cm} [唐] 李商隐
\end{center}

\begin{quote}
\large
云母屏风烛影深，\\
长河渐落晓星沉。\\
嫦娥应悔偷灵药，\\
碧海青天夜夜心。
\end{quote}

\vspace{0.3cm}

\begin{quote}
\textit{云母屏风上，烛影深深。银河渐渐落下，晨星也沉没了。嫦娥应该后悔偷了长生不老药吧——如今只有碧蓝的大海和辽阔的天空，日日夜夜陪伴着她孤独的心。}
\end{quote}

\vspace{0.3cm}

\textbf{背景小知识：}

在中国神话中，嫦娥是后羿的妻子。后羿从西王母那里得到了长生不老药，嫦娥偷吃了它，结果身体变轻，飞到了月亮上，从此一个人孤独地生活在月宫里。

\vspace{0.3cm}

\subsection*{\textcolor{teal!80!black}{古朗月行（节选）}}

\begin{center}
\large
\textbf{古朗月行（节选）} \hspace{2cm} [唐] 李白
\end{center}

\begin{quote}
\large
小时不识月，呼作白玉盘。\\
又疑瑶台镜，飞在青云端。\\
仙人垂两足，桂树何团团。\\
白兔捣药成，问言与谁餐？
\end{quote}

\vspace{0.3cm}

\begin{quote}
\textit{小时候不认识月亮，叫它"白玉盘"。又怀疑是瑶台的镜子，飞到了青云之上。月亮里仙人的双脚隐约可见，桂树长得圆圆的。月宫里的白兔在捣药，捣好了给谁吃呢？}
\end{quote}

\vspace{0.3cm}

\textbf{背景小知识：}

李白这首诗引用了好几个月亮的神话：月宫里有仙人、有桂树、有捣药的白兔——这些都是中国古代关于月亮的传说。

\vspace{0.5cm}

\thinktitle
\begin{enumerate}
  \item 李商隐说"嫦娥应悔偷灵药"——他怎么知道嫦娥后悔了？你觉得嫦娥会不会后悔？
  \item 李白把月亮比作"白玉盘"和"瑶台镜"。你觉得月亮还像什么？
  \item 同样写月亮，两首诗给你的感觉有什么不同？
\end{enumerate}

\inferbox
\vspace{0.1cm}
\noindent
\textbf{诗人为什么要在诗里写神话？他们明明可以直接写自己的心情，为什么要绕个弯子？你读诗的时候，有没有因为"看懂了这个典故"而觉得更开心？}

\genrebox
\vspace{0.1cm}
\noindent
\textbf{什么是"用典"？}

用典就是在自己的作品里借用古代的故事或诗句。比如李白写月亮时用了"神仙""桂树""白兔"这些典故，读过神话的人一看就知道他在写月亮。写作者用典故，就像和朋友说暗号——懂的人会心一笑。

\writetitle
\vspace{0.1cm}
选一个月亮的意象（如白玉盘、瑶台镜、桂树、白兔），写4句小诗，用上神话典故。

\newpage

% ============================================================
\section*{\textcolor{orange!80!black}{第四周：埃及神话——死亡与复活}}
% ============================================================

我们已经看了中国和希腊的神话，现在我们把目光转向古埃及。古埃及人最关心的一个问题是：人死后会去哪里？他们用神话来回答这个问题。

\vspace{0.5cm}

\subsection*{\textcolor{teal!80!black}{奥西里斯的故事（节选）}}

\begin{quote}
\textit{很久以前，埃及有一位英明的国王叫奥西里斯。他教会了埃及人耕种、酿酒、建造城市，人民非常爱戴他。}

\textit{然而，奥西里斯的弟弟塞特非常嫉妒他。塞特偷偷量了奥西里斯的身材，命人做了一个精美的箱子。在一次宴会上，塞特宣布：谁刚好能躺进这个箱子，箱子就送给谁。大家轮流试，都太大了或太小了。轮到奥西里斯时，他一躺进去——正好！塞特立刻冲上来盖紧箱盖，用钉子封死，把箱子扔进了尼罗河。}

\textit{奥西里斯的妻子伊西斯非常悲痛，她沿着河寻找丈夫。终于，她在一个遥远的地方找到了箱子。然而塞特发现了，他把奥西里斯的尸体切成了十四块，散落在埃及各处。}

\textit{伊西斯不肯放弃。她划着小船，一块一块地找回丈夫的尸体。每找到一块，她就埋葬在那里，并建一座神庙。最终，在神灵的帮助下，奥西里斯复活了——但不再是凡人，而是成为了冥界的国王，审判死者的灵魂。}
\end{quote}

\vspace{0.3cm}

\thinktitle
\begin{enumerate}
  \item 塞特为什么要害奥西里斯？这和前面哪些故事类似？
  \item 伊西斯为什么坚持要把丈夫的尸体找回来？这需要什么样的品质？
  \item 和中国神话、希腊神话相比，埃及神话最大的特点是什么？（提示：死亡主题）
\end{enumerate}

\inferbox
\vspace{0.1cm}
\noindent
\textbf{#1：塞特用"量身做箱"的诡计害了奥西里斯——你从中学到了什么？}

\vspace{0.1cm}
\noindent
\textbf{#2：奥西里斯的尸体被切成14块散落各地——为什么神话里要用这么"残酷"的细节？}

\genrebox
\vspace{0.1cm}
\noindent
\textbf{三大神话对比：}

\begin{tabular}{|p{0.28\textwidth}|p{0.28\textwidth}|p{0.28\textwidth}|}
\hline
\textbf{中国神话} & \textbf{希腊神话} & \textbf{埃及神话} \\
\hline
简短、含蓄 & 故事丰富、像连续剧 & 围绕死亡与复活 \\
人与自然的关系 & 神与人、神与神的关系 & 生与死的关系 \\
精卫/夸父靠自己的意志 & 神有人的缺点 & 复活依靠神灵帮助 \\
\hline
\end{tabular}

\writetitle
\vspace{0.1cm}
画一张"世界神话对照表"——找出中国神话（精卫/夸父）、希腊神话（普罗米修斯/潘多拉）、埃及神话（奥西里斯）在以下方面的异同：
\begin{itemize}
  \item 主人公的身份
  \item 主人公面对的困难
  \item 结尾是怎样的
  \item 想告诉我们什么
\end{itemize}

\newpage

% ============================================================
\section*{\textcolor{orange!80!black}{第五周：单元整合——创造你自己的神话}}
% ============================================================

这四周来，我们读了来自三个文明的神话故事。现在，让我们一起来整理一下、思考一下，然后用学到的"神话感"来创作一个属于自己的神话。

\vspace{0.5cm}

\subsection*{\textcolor{teal!80!black}{回顾与归纳}}

\textbf{一、神话的共同特点}

请你想一想：精卫填海、普罗米修斯盗火、奥西里斯复活……这些故事虽然来自不同的文明，但有没有什么共同之处？

\vspace{0.2cm}
\begin{enumerate}
  \item 它们都涉及\underline{\hspace{4cm}}的力量（凡人做不到的事）
  \item 它们都在解释某种\underline{\hspace{4cm}}（为什么有桃林？为什么人间有苦难？人死后去哪？）
  \item 它们的主人公都有\underline{\hspace{4cm}}的精神（不放弃、勇敢、牺牲）
  \item 它们最初都被人们\underline{\hspace{4cm}}（大家真的相信这些故事）
\end{enumerate}

\vspace{0.5cm}

\textbf{二、神话中的"对手"}

不同的神话中，主角面对不同形式的"对手"：

\begin{itemize}
  \item 精卫 vs \underline{\hspace{3cm}}（自然的力量）
  \item 夸父 vs \underline{\hspace{3cm}}（也是自然——太阳）
  \item 普罗米修斯 vs \underline{\hspace{3cm}}（权威——宙斯）
  \item 奥西里斯 vs \underline{\hspace{3cm}}（邪恶——塞特）
\end{itemize}

这些"对手"有什么区别？主角又是怎么应对的？

\vspace{0.5cm}

\subsection*{\textcolor{teal!80!black}{创作任务：我的神话}}

现在轮到你了！请你创作一个你自己的神话。

\textbf{步骤一：选择一个自然现象或事物起源来解释}

比如：
\begin{itemize}
  \item 为什么会有彩虹？
  \item 为什么天会下雨？
  \item 为什么人有影子？
  \item 为什么树叶会在秋天变黄？
  \item 为什么猫会抓老鼠？
\end{itemize}

\textbf{步骤二：构思你的主人公}

主人公是谁？是人？是动物？还是一个你想象出来的生灵？他/她有什么特别的能力？

\textbf{步骤三：设计一个困难或冲突}

主人公遇到了什么阻碍？谁是他的对手？

\textbf{步骤四：想一个结尾}

主人公成功了吗？失败了吗？如果是失败，有没有留下什么美好或有用的东西？（像精卫留下的精神、夸父留下的桃林）

\textbf{步骤五：写下来}

300-500字，试着用你在这几周学到的"神话感"来写。

\vspace{0.5cm}

\thinktitle
\vspace{0.1cm}
\begin{enumerate}
  \item 写完之后，把你的神话讲给爸爸妈妈听——看看他们能不能从中发现"神话的感觉"。
  \item 对比一下，你创造的这个神话和课堂上学到的哪个神话最像？为什么？
\end{enumerate}

\vspace{1cm}
\begin{center}
\textcolor{blue!70!black}{\large —— 神话就是人类最早的想象力。—— }
\end{center}

\end{document}
